% ============================================================================
% Technical appendix for national-AI-validation
% Requires: amsmath, amssymb, booktabs
% ============================================================================
\providecommand{\navdir}{national-AI-validation}

\section{Technical Appendix to the National Probability-Based Validation Simulation}
\label{app:sim-national-ai-validation}

\subsection{Finite-population generation}

The finite population contains \(H=400\) hospitals, of which
\(H_0=300\) belong to the standard design stratum and \(H_1=100\) to the
priority stratum. Hospital sizes are generated from a truncated log-normal
construction,
\[
  N_h
  =
  \min\{105,\max\{35,\operatorname{round}(\widetilde N_h)\}\},
  \qquad
  \log \widetilde N_h
  \sim
  \mathcal N(4.1276331,0.25^2).
\]
The realised total is \(N_T=25{,}585\).

For each hospital, let
\[
  (b_{Qh},b_{Yh},b_{Ah})^\top
  \sim
  \mathcal N(0,\Sigma_b),
\]
with standard deviations \((0.30,0.38,1.20)\) and correlation matrix
\[
  R_b
  =
  \begin{pmatrix}
    1 & 0.15 & 0 \\
    0.15 & 1 & 0.15 \\
    0 & 0.15 & 1
  \end{pmatrix}.
\]
For patient \(i\) in hospital \(h\), the harder-subgroup indicator is
\begin{equation}
  \operatorname{logit}\Pr(Q_{hi}=1\mid H_h,b_{Qh})
  =
  \operatorname{logit}(0.12)+1.80H_h+b_{Qh}.
  \label{eq:nav-app-q}
\end{equation}
The deterioration outcome is generated from
\begin{equation}
  \operatorname{logit}\Pr(Y_{hi}=1\mid H_h,Q_{hi},b_{Yh})
  =
  \operatorname{logit}(0.14)+0.45H_h+0.65Q_{hi}+b_{Yh}.
  \label{eq:nav-app-y}
\end{equation}
Among outcome-positive admissions, the locked alert is generated from
\begin{equation}
  \operatorname{logit}\Pr(A_{hi}=1\mid Y_{hi}=1,H_h,Q_{hi},b_{Ah})
  =
  \operatorname{logit}(0.97)-0.55H_h-2.00Q_{hi}+b_{Ah}.
  \label{eq:nav-app-a}
\end{equation}
For completeness, false-alert probabilities among \(Y=0\) are generated from
\[
  \operatorname{logit}\Pr(A_{hi}=1\mid Y_{hi}=0,H_h,Q_{hi},b_{Yh})
  =
  \operatorname{logit}(0.10)+0.15H_h+0.15Q_{hi}+0.20b_{Yh}.
\]
False alerts do not enter the sensitivity estimands.

The population is generated once with master seed 141421 and then frozen. No
DGP parameter is altered after observing the strategy results.

\subsection{Target and observation parameters}

The finite-population target distribution assigns mass \(1/N_T\) to every
eligible admission. The national and subgroup parameters are ratios of fixed
finite-population totals:
\[
  \theta_{T,q}
  =
  \frac{T_{T,q}}{Y_{T,q}},
  \qquad
  T_{T,q}=\sum_h\sum_i\mathbf 1(Q_{hi}=q)A_{hi}Y_{hi},
  \qquad
  Y_{T,q}=\sum_h\sum_i\mathbf 1(Q_{hi}=q)Y_{hi}.
\]
The national parameter is obtained by omitting the subgroup indicator.

For design \(k\), the unweighted observation distribution is
\[
  P_{\mathcal O,k}(h,i)
  =
  \frac{\pi_{hi}^{(k)}}{\sum_g\sum_r\pi_{gr}^{(k)}}.
\]
Its subgroup sensitivity parameter is
\[
  \theta_{\mathcal O,k,q}
  =
  \frac{
    \sum_h\sum_i \pi_{hi}^{(k)}\mathbf 1(Q_{hi}=q)A_{hi}Y_{hi}
  }{
    \sum_h\sum_i \pi_{hi}^{(k)}\mathbf 1(Q_{hi}=q)Y_{hi}
  }.
\]
These parameters were calculated exactly from the locked finite population
before Monte Carlo sampling.

\subsection{Inclusion probabilities and allocations}

N1 draws a patient SRSWOR of size
\[
  n_m
  =
  \operatorname{round}\left(N_T\frac{m}{400}\right),
\]
which gives \(n_m\in\{2559,5117,7676,10234\}\) for
\(m\in\{40,80,120,160\}\). Every patient has inclusion probability
\(n_m/N_T\).

For hospital strategies, hospitals are sampled independently by SRSWOR within
stratum. If \(m_j\) hospitals are selected from \(H_j\), then every admission
in hospital stratum \(j\) has first-order inclusion probability
\[
  \pi_{hi}=\pi_h=\frac{m_j}{H_j}.
\]
N2 uses proportional allocations
\[
  (m_0,m_1)
  \in
  \{(30,10),(60,20),(90,30),(120,40)\}.
\]
N3--N5 use enriched allocations
\[
  (m_0,m_1)
  \in
  \{(20,20),(40,40),(60,60),(80,80)\}.
\]
All admissions in selected hospitals are observed during the same fixed audit
period. There is no patient subsampling inside hospitals and no month-level
selection stage.

\subsection{Estimators}

For N1, ordinary sample sensitivity is
\[
  \widehat\theta_{N1}
  =
  \frac{\sum_{i\in s}A_iY_i}{\sum_{i\in s}Y_i}.
\]

For N3, the enriched ORP is analysed using the pooled unweighted ratio
\[
  \widehat\theta_{N3}
  =
  \frac{\sum_{h\in s_0\cup s_1}T_h}
       {\sum_{h\in s_0\cup s_1}Y_h}.
\]

For N2, N4, and N5, define design-expanded stratum totals
\[
  \widehat T_j
  =
  \frac{H_j}{m_j}\sum_{h\in s_j}T_h,
  \qquad
  \widehat Y_j
  =
  \frac{H_j}{m_j}\sum_{h\in s_j}Y_h.
\]
The design-weighted combined ratio is
\[
  \widehat\theta_w
  =
  \frac{\widehat T_0+\widehat T_1}
       {\widehat Y_0+\widehat Y_1}.
\]
For subgroup \(q\), \(T_h\) and \(Y_h\) are replaced by
\(T_{hq}\) and \(Y_{hq}\) before stratum totals are added.

\subsection{Variance procedures}

For the design-weighted combined ratio, let
\[
  z_{hj}=T_{hj}-\widehat\theta_wY_{hj}
\]
and let \(s_{z,j}^2\) be the sample variance of \(z_{hj}\) among selected
hospitals in stratum \(j\). The Taylor variance estimator is
\begin{equation}
  \widehat{\operatorname{Var}}(\widehat\theta_w)
  =
  \frac{
    \sum_{j=0}^1
    H_j^2\left(1-\frac{m_j}{H_j}\right)
    \frac{s_{z,j}^2}{m_j}
  }{
    (\widehat Y_0+\widehat Y_1)^2
  }.
  \label{eq:nav-app-taylor}
\end{equation}
Independent stratum contributions are therefore added before division by the
squared expanded denominator.

For N3, the same ratio-linearisation idea is applied to unexpanded sample
sums. Its variance estimates sampling dispersion around the parameter targeted
by the unweighted estimator:
\[
  \widehat{\operatorname{Var}}(\widehat\theta_{N3})
  =
  \frac{
    \sum_{j=0}^1
    m_j\left(1-\frac{m_j}{H_j}\right)s_{z,j}^2
  }{
    \left(\sum_{h\in s_0\cup s_1}Y_h\right)^2
  }.
\]

N5 uses the same weighted point estimate as N4 but treats positive patients as
independent. With stratum weights \(w_j=H_j/m_j\), its naive variance is
\[
  \widehat{\operatorname{Var}}_{\mathrm{naive}}(\widehat\theta_w)
  =
  \frac{
    \sum_{j=0}^1 w_j^2
    \left[
      T_{s,j}(1-\widehat\theta_w)^2
      +(Y_{s,j}-T_{s,j})\widehat\theta_w^2
    \right]
  }{
    \left(\sum_{j=0}^1w_jY_{s,j}\right)^2
  }.
\]
This expression ignores the hospital-level dependence and is intentionally
incompatible with the design.

For all strategies, a logit-transformed 95\% interval is formed as
\[
  \operatorname{expit}\left[
    \operatorname{logit}(\widehat\theta)
    \pm
    1.96\frac{\widehat{\operatorname{SE}}(\widehat\theta)}
                    {\widehat\theta(1-\widehat\theta)}
  \right].
\]

\subsection{PR, uncertainty, and TAC classification}

For each strategy and estimand, the Monte Carlo bias and variance are
\[
  \operatorname{Bias}_k
  =
  \frac{1}{R}\sum_{r=1}^R\widehat\theta_{kr}-\theta_T,
  \qquad
  \operatorname{Var}_k
  =
  \frac{1}{R-1}\sum_{r=1}^R
  (\widehat\theta_{kr}-\overline{\widehat\theta}_k)^2.
\]
The declared point-estimation PR conditions are
\[
  |\operatorname{Bias}_k|\leq0.03,
  \qquad
  \operatorname{Var}_k\leq(0.03/1.96)^2.
\]
Interval coverage is reported separately. Formal TAC requires the PR
point-estimation conditions, acceptable target-parameter coverage, and a
variance procedure structurally compatible with the design. The last condition
prevents accidental coverage in one simulation cell from legitimising N5's
known incompatible variance rule.

For a valid interval \([L,U]\),
\[
  \operatorname{TAC}(0.85)
  =
  \begin{cases}
    \text{Adequate}, & L\geq0.85,\\
    \text{Not adequate}, & U<0.85,\\
    \text{Inconclusive}, & L<0.85\leq U.
  \end{cases}
\]
When the evidential gate fails, the formal status is
\(\text{Evidentially insufficient}\). Mechanical interval-threshold outcomes
are retained in the replication file for diagnostic purposes only.

\subsection{Monte Carlo uncertainty}

For \(R\) estimable replications, the Monte Carlo standard error of bias is
\[
  \operatorname{MCSE}(\widehat{\operatorname{Bias}})
  =
  \frac{\operatorname{ESD}(\widehat\theta)}{\sqrt R}.
\]
For an estimated proportion \(\widehat p\), including coverage, TAC, and
non-estimability frequencies,
\[
  \operatorname{MCSE}(\widehat p)
  =
  \sqrt{\frac{\widehat p(1-\widehat p)}{R}}.
\]
Writing \(e_r^2=(\widehat\theta_r-\theta_T)^2\), the delta-method MCSE of RMSE
is
\[
  \operatorname{MCSE}(\widehat{\operatorname{RMSE}})
  =
  \frac{\operatorname{SD}(e_r^2)}
       {2\widehat{\operatorname{RMSE}}\sqrt R}.
\]

\subsection{Seed mapping and implementation checks}

The Python sub-seed for master seed \(M\), design label \(L\), sample size
\(m\), replication \(r\), and language tag is generated from
\[
  \texttt{SHA256}(M\,|\,L\,|\,m\,|\,r\,|\,\texttt{python}).
\]
The first eight digest bytes are interpreted as an unsigned little-endian
integer and reduced modulo \(2^{63}-1\). N3--N5 share the label
\texttt{ENRICHED}, ensuring identical samples within replication. The R
implementation uses the analogous SHA-256 mapping with language tag
\texttt{R} and an R-compatible integer reduction.

The executable validation suite verifies inclusion probabilities, absence of
month sampling, exact sample identity for N3--N5, exact point-estimate identity
for N4 and N5, direct enumeration of finite-population truths, approximate
unbiasedness of N1, design consistency of N4, convergence of N3 toward
\(\theta_{\mathcal O,3}\), detection of N5 uncertainty failure, and the absence
of formal TAC classifications after an evidential-gate failure. The full
machine-readable results and software checksums are stored in
\texttt{results/summary/validation\_results.json}.
